\chapter{A collective trajectory is not a complete interface}
\label{comp:public-compression}

Everything in this chapter is a statement about the reciprocal law potential, the cooperative consensus potential whose Fisher flow moves each source to reduce the costs of the receivers that attend to it. It is not a statement about the free energy of any agent's generative model. Within that potential there is a way to compress distinct public laws under the smooth interaction. It uses their collective motion. Near a common stable reference, differences between agents decay faster than their shared evolution. One can then assign each fine configuration a virtual collective state whose future follows exactly the same reduced trajectory. The fine agents need not coincide at any finite time, but they all approach the reference, so what is compressed is the decaying part of the approach. The construction is the asymptotic phase familiar from the theory of isochrons and isostables, carried out in charts of bounded log densities.

For an explicit general-law construction, take probability references $A,B$ with balanced binary features $\varphi,\psi$, and the fixed normalized evaluator
\begin{equation}
 G(m,dk)=\frac{e^{J\varphi(k)\psi(m)}}{\cosh J}A(dk),\qquad |J|<1.
 \label{comp:eq-phase-evaluator}
\end{equation}
The bound $|J|<1$ is exactly the condition for the common reference to be stable. The other sample coordinates can be arbitrary. Write $q_i=e^{u_i}A/\E_Ae^{u_i}$ and $s_i=e^{v_i}B/\E_Be^{v_i}$, with bounded reference-centered log densities. No finite-dimensional family is imposed. Use identity channels, uniform nonself source priors, and the Gibbs-reduced reciprocal potential.

Set $z_i=(u_i,v_i)$. On the virtual diagonal the common law follows
\begin{equation}
 \dot u=-u+J\varphi\E_{s_v}\psi,\qquad
 \dot v=-v+J\psi\E_{q_u}\varphi.
 \label{comp:eq-phase-common}
\end{equation}
At the reference state, the linearized fine dynamics is
\begin{equation}
 \dot z_i=\mathcal A z_i-\frac{2N}{N-1}(z_i-\bar z),\qquad
 \bar z=\frac1N\sum_i z_i.
 \label{comp:eq-phase-rates}
\end{equation}
Here $\mathcal A$ is the linearization of \cref{comp:eq-phase-common}. Its slowest common decay rate is $1-|J|$, whereas the slowest difference rate is $1+2N/(N-1)-|J|$. The factor two in the peer term includes both receiving and being used as a source. This gap is what the interaction contributes. If the peer costs are switched off, each agent relaxes on its own with the common rates, the differences decay no faster than the average, and the construction below loses its foundation.

The arithmetic average of the log densities is not itself an exact closed state. Away from the diagonal, its equation contains a correction quadratic in the child differences. That correction decays fast enough to define a nonlinear phase. Let $\Phi_t$ be the fine flow and $\phi_t$ the flow of \cref{comp:eq-phase-common}. The exact parent coordinate is
\begin{equation}
 \Pi(x)=\lim_{t\to\infty}\phi_{-t}\bigl(\overline{\Phi_t x}\bigr),\qquad
 \Pi(\Phi_t x)=\phi_t(\Pi x).
 \label{comp:eq-phase-map}
\end{equation}
The backward common flow transports the late collective state back to the initial time. The fast decay of child differences makes the limit converge, and the limit is the same whichever smooth symmetric average is used. Its local fibers contain continuously many distinct child configurations, so this construction discards public degrees of freedom rather than merely their labels. The parent evolves exactly once its initial phase is known. The phase can be computed to exponential accuracy by integrating forward for a time $t_0$, averaging, and integrating the common flow backward for the same time. It is also the only choice: any Lipschitz parent map that is the identity on equal children and intertwines the two flows equals $\Pi$.

The response of the group is a separate question. For two binary agents with equal belief and model means within each agent, put $x_1=m+a$ and $x_2=m-a$. Their phase-law mean is
\begin{equation}
 \bar p=m+\frac{3(2-J)}{2(5-J)}ma^2+O(|m|a^2(|m|+|a|)^2).
 \label{comp:eq-phase-response}
\end{equation}
A uniform choice of child followed by one bit sample has mean $m$. The diagonal configuration with both child means equal to $\bar p$ has the same phase parent, but its corresponding response has mean $\bar p$. For small $m=a>0$, these means differ. The parent predicts its own inherited trajectory while failing to reproduce this ordinary group response. Sampling the phase law would change the interface. Because the phase is unique, no exact single-law parent of this kind can do better, although the defect is small: it is second order in the child differences.

An unrestricted identity-channel Gibbs receiver can distinguish even more. For child densities $f_j$ and positive weights $w_j$, its block normalizer is
\begin{equation}
 Z(r)=\sum_jw_j e^{-\KL(r\lambda\Vert f_j\lambda)/T}.
 \label{comp:eq-phase-gibbs}
\end{equation}
Knowing this response for all nearby receiver laws determines every distinct exposed child density and its total weight. To see the mechanism, cancel the known receiver entropy. The remaining function is a finite sum of exponentials of the linear probes $\int r\log f_j$. A direction that separates these probes gives a Vandermonde system, which recovers their coefficients. Distinct child laws therefore cannot be forgotten while retaining this entire interface. A fixed comparison channel changes the conclusion in a precise way: the receiver then sees only the channelled laws, and children with equal channelled laws are indistinguishable to it. For a deterministic coarse-graining channel even the force that the receiver exerts on its sources depends only on the channelled laws.

Restricting the possible receivers may permit compression. It must be accompanied by an exact check of the incoming forces and of repeated-response memory. For a second formation step, joining parents must also respect the information already discarded. Measured against the formation standard, the collective phase is a retention of a supplied block rather than a formation: nothing is selected, the ordinary sampling response is not inherited, and the fine differences it compresses decay. It supplies a continuous autonomous reduction without an erasure gate, but not the recursive interface properties or a dynamically selected hierarchy.
